\documentclass{book}
\usepackage[utf8]{inputenc}
\usepackage[T1]{fontenc}
\usepackage{geometry}
\usepackage{graphicx}
\usepackage{float}
\usepackage{hyperref}
\usepackage[english,indonesian]{babel}

\title{Optimasi Pemrosesan Template LaTeX Menggunakan Pendekatan Berbasis AI}
\author{Academic Support Agent}
\date{\today}

\begin{document}
\maketitle

\begin{abstract}

\end{abstract}

\chapter{PENDAHULUAN}


\chapter{TINJAUAN PUSTAKA}


% Original content (commented out):
% \documentclass[11pt, a4paper, oneside]{book}

\input{prelim/preamble}

%%%%%%%%%%%%%%%%%%%%%%%%%%%%%%%%%%%%%%%%%%%%%%%%%%%%%%%%%%%%%%%%%%%%%%%%%%%%
\begin{document}

% SET THE FIRST PART TO BE SINGLE SPACED
{\setstretch{1.0}
\input{prelim/prelim}
}

\mainmatter

%%%%%%%%%%%%%%%%%%%%%%%%%%%%%%%%%%%%%%%%%%%%%%%%%%%%%%%%%%%%%%%%%%%%%%%%%%%%
\include{chapters/01_intro}
\cleartooddpage

% insert more chapters here

% glossary and acronyms managed in prelim/glossary, to print:
% 1. run latex command
% 2. use command: makeglossaries %
% 3. use normal build procedure (PDFLaTeX, BibTeX, PDFLaTeX, PDFLaTeX)
\printglossaries
\cleartooddpage

%%%%%%%%%%%%%%%%%%%%%%%%%%%%%%%%%%%%%%%%%%%%%%%%%%%%%%%%%%%%%%%%%%%%%%%%%%%%
\begin{appendices}
\input{appendices/appendices}
\end{appendices}
\cleartooddpage

%%%%%%%%%%%%%%%%%%%%%%%%%%%%%%%%%%%%%%%%%%%%%%%%%%%%%%%%%%%%%%%%%%%%%%%%%%%%
{\setstretch{1.0}
\bibliographystyle{prelim/apa-good}
\bibpunct[ ]{(}{)}{;}{a}{}{,}
\phantomsection
\addcontentsline{toc}{chapter}{Bibliography}
\bibliography{BibFile}
\cleartooddpage
}

%%%%%%%%%%%%%%%%%%%%%%%%%%%%%%%%%%%%%%%%%%%%%%%%%%%%%%%%%%%%%%%%%%%%%%%%%%%%
\backmatter

\end{document}

\end{document}